\section{Appendix Q -- Synthesis Support Digest}
% R005_APPENDIX_READER_FRAME
Reader frame. This appendix supports the main manuscript by explaining source role, evidence class, audit route, and limitation boundary. It is not a detached raw dump.


\label{sec:oc-core-1-3-synthesis-support-dossiers}

This digest explains how to inspect the bounded K-level synthesis without importing raw curated evidence dossiers into the public monograph. The main synthesis chapter states the scientific claim. The evidence package carries the full machine-readable support rows and checksums.


A reader should inspect this route by asking whether a K-level row names its role, proof route, replay boundary, comparator boundary, and reopening condition. If a raw dossier row is needed, the corpus register points to it outside the reading path so the PDF remains a coherent manuscript.


\subsection{Scholarly integration essay}

The synthesis support material has a simple editorial purpose: it prevents a long cross-domain model from becoming a sequence of disconnected demonstrations. The reader should see the same discipline at every scale. A formal claim names the object grammar and its assumptions. A proof claim names the theorem boundary. A computational claim names the finite or executable witness. An empirical claim names the observable, comparator, uncertainty, negative control, and falsifier. A prior-art claim names the overlap and the bounded residual difference. These are not separate habits; together they are the reading grammar of the release.


The first integration burden is conceptual. OC Core 1.3.3 uses typed carriers, realizations, lawful possibility, time-indexed liveness, residue, morphism, boundary, operator, cycle, dimension, and K-level vocabulary. The support material is useful only if it makes those terms harder to misuse. A reader should be able to ask, at any point in the monograph, which typed object is under discussion, which relation preserves it, which boundary separates it, and which condition would make the statement false or overbroad.


The second integration burden is evidential. The release is ambitious precisely because it refuses to let ambition float free from evidence. The finite semantic checks do not prove every empirical domain. The target-blind and replay rows do not mechanize every theorem. The comparator register does not establish global priority. Each evidence family has its own authority and its own ceiling. The scientific strength comes from making those ceilings visible and from refusing to promote a sentence when the corresponding evidence family cannot carry it.


The third integration burden is didactic. A hostile reader should not have to guess why an appendix exists. The appendix should lower the cost of verification: it should tell the reader which kind of question to ask, which artifact class to open, and which failure would reopen the public wording. That is why long machine rows are kept in the evidence package while the manuscript explains their scientific role. The monograph is the argument; the package is the inspection apparatus.


The fourth integration burden is comparative. OC is not presented as a vocabulary that erases systems theory, autopoiesis, dynamical systems, category-style formalisms, RAF closure, complexity measures, identity theories, or reproducible-research standards. It is presented as a typed model-core release that packages selected claims with proof boundaries, executable witnesses, replay discipline, and explicit reopening rules. Where a predecessor already carries the same claim under the same assumptions, OC receives no priority credit for that sentence.


The fifth integration burden is lifecycle clarity. The reader should not confuse death, residue, rebirth, and identity. A dead continuum is not live by vocabulary alone. A residue is not the same token as a continuing live identity. A rebirth relation is evidence-bound and target-bound. An identity claim must preserve the declared invariant. These distinctions are central because they make the model criticizable: a reviewer can attack the invariant, the endpoint relation, the admissibility condition, or the evidence row instead of arguing against a vague metaphysical label.


The sixth integration burden is empirical humility without empirical retreat. The public replay rows are included because they force the package to demonstrate source discipline, numeric reconstruction, comparators, residuals, controls, and falsifiers. They are not included to pretend that all science has already been numerically exhausted. This is a stronger public posture than overclaiming: it gives a critic exact handles, and it gives future research a precise path for promotion.


The seventh integration burden is release accountability. A scientific release can fail even when its code compiles if the public PDF is unreadable, the title page is absent, the bibliography is thin, the theorem boundary is hidden, the evidence rows are printed without interpretation, or the archive metadata misleads the reader. The synthesis support material therefore belongs to the manuscript only where it improves the reader's understanding. Everything else belongs in the evidence package.


The final integration burden is continuity from this release to later work. OC Core 1.3.3 is a bounded external-review release, while the larger full-scope science program remains active. This distinction does not weaken the release. It protects it. A later release may promote stronger claims only by adding the missing proof, replay, comparator, and falsifier layers. Until then, the public monograph must show both the current achievement and the exact boundary of what has not yet been earned.


A practical reading exercise makes this concrete. Take any strong sentence in the release and ask four questions in order. First, what is the typed object or relation named by the sentence? Second, which proof, finite witness, replay row, or comparator row is allowed to support it? Third, what is the strongest adjacent sentence that the evidence does not support? Fourth, what observation, countermodel, failed build, failed replay, or stronger comparator would reopen it? If the reader can answer all four, the sentence belongs in the release. If not, it belongs in the research backlog or evidence package.


This exercise also explains why the manuscript and the machine-readable package have different jobs. The manuscript teaches the scientific object and the boundaries of its current evidence. The package lets a reviewer inspect exact rows, hashes, and files. Confusing the two was the criticism route corrected in the 1.3.3 repair: rows and hashes are necessary for reproducibility, but they do not become a readable scientific argument until the manuscript explains their conceptual role, their evidential ceiling, and their reopening condition.


The editorial consequence is that every major document in the release must answer the same set of reader questions from a different distance. The guide answers what to open first and why. The journal core answers what the central contribution is. The monograph answers how the model, proof surface, evidence surface, comparison boundary, and reviewer response cohere. The methods companion answers how to reproduce or challenge the evidence. The reviewer map answers where the strongest objections land. A defect in any one document is therefore not a cosmetic issue; it weakens the whole public scientific surface unless the other documents still give a coherent route to the same claim boundary.


A second practical exercise is to read the release backward. Start with a hostile objection, move to the attacked claim, then to the proof sheet or replay row, then to the formal definition that makes the claim meaningful. If the backward path breaks, the objection has found a real weakness. If the path holds, the reader can see exactly why the release is bounded but still strong. This backward reading is especially important for claims about liveness, residue, rebirth, K-level witnesses, operator semantics, and empirical replay, because each of those areas is easy to overstate when presented only as narrative.


The support synthesis also protects future work from drift. When a later release expands empirical lanes or promotes a stronger theorem, the new claim should enter through the same gates: typed object, assumption set, proof or executable witness, comparator boundary, negative control where empirical, and falsifier. This means that the 1.3.3 monograph is not merely a document about the current model; it is a public example of how OC claims must be promoted if they are to remain scientifically inspectable.


\subsection{Inspection route}

Start with the current T133 proof surface, then compare the relevant K-level role with the finite semantic case and lawful-demotion rule. Use the public comparator rows only for positioning, not for priority. Treat future full-domain synthesis as a background research obligation unless a later release supplies theorem, replay, comparator, negative-control, and falsifier evidence.


\subsection{K-level synthesis source digest}

The full synthesis support route is represented here as a reading digest. The exact source and checksum remain in the evidence package; this text explains how a reviewer should use the source without turning the monograph into a generated dossier.
